\begin{agradecimentos}
    \noindent
    Este livro é fruto de anos de pesquisa, experimentação e desenvolvimento 
    contínuo do \opencodeeco{}. Agradeço profundamente à comunidade de 
    engenharia de software e inteligência artificial que, através de décadas 
    de pesquisa aberta e colaborativa, tornou possível a construção de 
    sistemas cognitivos artificiais cada vez mais sofisticados.
    
    Agradecimentos especiais aos pesquisadores e engenheiros cujo trabalho 
    fundamenta cada capítulo desta obra: John McCarthy, Marvin Minsky, 
    Alan Turing, Claude Shannon, John von Neumann, Donald Knuth, Kent Beck, 
    Martin Fowler, Stuart Russell, Peter Norvig, Yoshua Bengio, Geoffrey 
    Hinton, Yann LeCun, Andrew Ng, e tantos outros que pavimentaram o 
    caminho da inteligência artificial e da engenharia de software.
    
    À Universidade Federal do Ceará (UFC), campus Crateús, pelo ambiente 
    acadêmico que estimula a pesquisa interdisciplinar. Ao Programa de 
    Pós-Graduação em Tecnologia e Engenharia (PPGTE), pela estrutura 
    que permitiu a validação científica dos conceitos aqui apresentados.
    
    Aos estudantes e pesquisadores que testaram, questionaram e aprimoraram 
    cada componente do \opencodeeco{} ao longo de 23 ciclos evolutivos 
    (R1 a R23), meu mais sincero reconhecimento.
    
    Finalmente, à minha família, pelo apoio incondicional durante os 
    longos períodos de imersão necessários à conclusão desta obra.
\end{agradecimentos}
